\scanpage{131}

De plus, procédant par récurrence sur $d+d'$, et utilisant une \uline{factorisation d'Artin} de $U$ [ce qui est loisible, car on peut prendre $U$ petit], on est ramené au cas où $U$ est de \uline{dim 1}.

Alors la suite exacte pour $\deleted{\unsure{g}}fg$, \deleted{\unsure{se}} compte tenu de $R^0_!f(R^{\unsure{q|?}} g_*(\mathbb Q_\ell))=0$, se réduit à \deleted{\unsure{deux}} suites exactes courtes

\[
\begin{aligned}
0\mathrel{\text{---}}R^2_!f_*(R^{n-2}g_*(\mathbb Q_\ell))&\mathrel{\text{---}}\unsure{R^n|R^-|R^n_!}\unsure{(gf)_*|(gf)_{-}}(\mathbb Q_\ell)\mathrel{\text{---}}R^1_! f_*(R^{n-1}g_*(\mathbb Q_\ell))\\[-2pt]
&\hspace{17em}\downarrow\\[-2pt]&\hspace{17em}0
\end{aligned}
\]

donc par lemme 1 on est ramené à vérifier que $R^{\unsure{n|-}}_!\overset{\inserted{h}}{(gf)}_{\unsure{*|!}}(\mathbb Q_\ell)$ est spécial. Or utilisant la résolution des singularités pour $V$, on a $V=\widehat V-Z$, $\widehat V$ projectif et lisse sur $k$, d'où

\[ \mathbb R^{n-1}_! h_Z(\mathbb Q_\ell)\longrightarrow R^n_!\widehat h(\mathbb Q_\ell)\longrightarrow R^{\unsure{n|-}}_!\widehat h(\mathbb Q_\ell) \]

% archive sheet number: 130
